\section{Equivalence}

% equivalence/index.tex

\begin{toolkitbox}{Equivalence and Partitions --- Quick Reference}
\small
\begin{tabular}{l l}
\toprule
\textbf{Name}  & \textbf{Meaning} \\
\midrule
\hyperref[def:equiv-class]{Equivalence class}   & $[a]_R = \{x \in A \mid (a,x)\in R\}$                        \\
\hyperref[def:quotient-set]{Quotient set}        & $A/R$: all equivalence classes                                \\
\hyperref[def:index]{Index of $R$}        & Cardinality of $A/R$                                          \\
\hyperref[def:canonical-surj]{Canonical surjection}  & $\pi : A \to A/R$, $\pi(a) = [a]$                          \\
\hyperref[def:partition]{Partition}           & Nonempty, disjoint, covering collection of subsets            \\
\hyperref[lem:rep-independence]{Rep.\ independence}  & $[a]=[b] \iff (a,b)\in R$                                    \\
\hyperref[thm:equiv-partition]{Equiv.--partition correspondence}  & Bijection between equiv.\ relations and partitions \\
\bottomrule
\end{tabular}
\end{toolkitbox}

% =========================================================
% Equivalence Classes and Partitions
% =========================================================

\section{Equivalence Classes and Partitions}

% ---------------------------------------------------------
% TOOLKIT
% ---------------------------------------------------------

\begin{definitionbox}{Definition (Equivalence Class)}
\begin{definition}[Equivalence Class]\label{def:equiv-class}
Let $R$ be an equivalence relation on $A$. For $a \in A$, the
\emph{equivalence class} of $a$ is:
\[
[a]_R \;:=\; \{\, x \in A \mid (a,x) \in R \,\}.
\]
When $R$ is clear from context, we write $[a]$.
\end{definition}
\end{definitionbox}

\begin{dependencies}
\begin{itemize}
  \item \hyperref[def:equivalence-rel]{Equivalence Relation}
  \item \hyperref[def:relation]{Relation}
\end{itemize}
\end{dependencies}

\begin{remark*}[English reading]
$[a]_R$ is the set of all elements that $R$ declares ``the same as $a$.''
Two elements lie in the same class iff they are related: this is precisely
the Representative Independence Lemma below.
\end{remark*}

\begin{definitionbox}{Definition (Quotient Set)}
\begin{definition}[Quotient Set]\label{def:quotient-set}
The \emph{quotient set} of $A$ by $R$ is the set of all equivalence classes:
\[
A / R \;:=\; \{\, [a]_R \mid a \in A \,\}.
\]
\end{definition}
\end{definitionbox}

\begin{dependencies}
\begin{itemize}
  \item \hyperref[def:equivalence-rel]{Equivalence Relation}
  \item \hyperref[def:equiv-class]{Equivalence Class}
\end{itemize}
\end{dependencies}

\begin{definitionbox}{Definition (Index of an Equivalence Relation)}
\begin{definition}[Index of an Equivalence Relation]\label{def:index}
The \emph{index} of $R$ on $A$ is the cardinality $|A/R|$, i.e.\ the number
of equivalence classes.
\end{definition}
\end{definitionbox}

\begin{remark*}[Standard quantified statement]
$\mathrm{index}(R) = |A/R| = |\{[a]_R \mid a \in A\}|$.
\end{remark*}

\begin{dependencies}
\begin{itemize}
  \item \hyperref[def:quotient-set]{Quotient Set}
  \item \hyperref[def:same-cardinality]{Same Cardinality}
\end{itemize}
\end{dependencies}

\begin{definitionbox}{Definition (Canonical Surjection)}
\begin{definition}[Canonical Surjection]\label{def:canonical-surj}
The \emph{canonical surjection} (quotient map) is the function
\[
\pi : A \to A/R,
\qquad
\pi(a) := [a].
\]
\end{definition}
\end{definitionbox}

\begin{remark*}[Forward reference]
This definition uses the function vocabulary developed in the next chapter.
It is included here because quotient sets are naturally introduced with their
canonical map; proofs using its function-theoretic properties may be revisited
after the function chapter.
\end{remark*}

\begin{dependencies}
\begin{itemize}
  \item \hyperref[def:function]{Function}
  \item \hyperref[def:quotient-set]{Quotient Set}
  \item \hyperref[def:equiv-class]{Equivalence Class}
  \item \hyperref[def:surjective]{Surjective Function}
\end{itemize}
\end{dependencies}

\begin{remark*}[Properties of $\pi$]
$\pi$ is surjective by construction. Elements $a, b \in A$ satisfy
$\pi(a) = \pi(b)$ iff $(a,b) \in R$. The canonical surjection is the
prototype for all quotient constructions: it reappears as the quotient
homomorphism in algebra and the quotient map in topology.
\end{remark*}

\begin{definitionbox}{Definition (Partition)}
\begin{definition}[Partition]\label{def:partition}
A \emph{partition} of $A$ is a collection $\mathcal{P}$ of subsets of $A$ such
that:
\begin{enumerate}[label=(\roman*)]
\item every block is nonempty: $\forall P \in \mathcal{P},\; P \neq \varnothing$;
\item distinct blocks are disjoint:
$\forall P,Q \in \mathcal{P},\; P \neq Q \rightarrow P \cap Q = \varnothing$;
\item the blocks cover $A$: $\bigcup_{P \in \mathcal{P}} P = A$.
\end{enumerate}
The sets $P \in \mathcal{P}$ are called the \emph{blocks} of the partition.
\end{definition}
\end{definitionbox}

\begin{dependencies}
\begin{itemize}
  \item \hyperref[def:subset]{Subset}
  \item \hyperref[def:empty-set]{Empty Set}
  \item \hyperref[def:intersection]{Intersection}
  \item \hyperref[def:union]{Union}
  \item \hyperref[def:cover-full]{Cover}
\end{itemize}
\end{dependencies}

\begin{remark*}[Partition vs.\ cover]
Every partition of $A$ is a cover of $A$ whose members are nonempty and
pairwise disjoint. Partitions are precisely the covers satisfying the
disjointness condition.
\end{remark*}

% ---------------------------------------------------------
% Core Lemma and Theorem
% ---------------------------------------------------------

\begin{lemmabox}{Lemma (Representative Independence Lemma)}
\begin{lemma}[Representative Independence Lemma]\label{lem:rep-independence}
Let $R$ be an equivalence relation on $A$. For any $a,b \in A$,
\[
[a] = [b]
\;\;\Longleftrightarrow\;\;
(a,b) \in R.
\]
\hyperref[prf:rep-independence]{Go to proof.}
\end{lemma}
\end{lemmabox}

\begin{remark*}[Standard quantified statement]
$\forall a, b \in A:\; [a]_R = [b]_R \;\leftrightarrow\; (a,b)\in R$.
\end{remark*}
\begin{dependencies}
\begin{itemize}
  \item \hyperref[def:equivalence-rel]{Equivalence Relation}
  \item \hyperref[def:equiv-class]{Equivalence Class}
  \item \hyperref[def:partition]{Partition}
  \item \hyperref[def:quotient-set]{Quotient Set}
\end{itemize}
\end{dependencies}

\begin{remark*}[Well-definedness on quotient sets]
This lemma is the key fact underlying well-definedness of functions on quotient
sets: a function $f : A \to B$ defined by $f([a]) = \cdots$ is well-defined
iff the formula gives the same output for all representatives of $[a]$.
\end{remark*}

\begin{theorembox}{Theorem (Equivalence Relations and Partitions)}
\begin{theorem}[Equivalence Relations and Partitions]\label{thm:equiv-partition}
Let $A$ be a set.
\begin{enumerate}[label=(\roman*)]
\item If $R$ is an equivalence relation on $A$, then $A/R$ is a partition of $A$.
\item If $\mathcal{P}$ is a partition of $A$, then the relation
$R_{\mathcal{P}}$ defined by
\[
(a,b) \in R_{\mathcal{P}} \;\Longleftrightarrow\;
\exists P \in \mathcal{P} \text{ with } a \in P \text{ and } b \in P
\]
is an equivalence relation on $A$.
\item These constructions are inverse: $R_{A/R} = R$ and $A/R_{\mathcal{P}} = \mathcal{P}$.
\end{enumerate}
\hyperref[prf:equiv-partition]{Go to proof.}
\end{theorem}
\end{theorembox}

\begin{remark*}[Standard quantified statement]
(i) $\forall a \in A,\; [a]_R \neq \varnothing$; $\forall a,b \in A,\; [a]_R \neq [b]_R \to [a]_R \cap [b]_R = \varnothing$; $\bigcup_{a \in A}[a]_R = A$.\\
(ii) $\forall a,b \in A:\; (a,b) \in R_{\mathcal{P}} \;\leftrightarrow\; \exists P \in \mathcal{P},\; a \in P \wedge b \in P$.
\end{remark*}
\begin{dependencies}
\begin{itemize}
  \item \hyperref[def:equivalence-rel]{Equivalence Relation}
  \item \hyperref[def:equiv-class]{Equivalence Class}
  \item \hyperref[def:quotient-set]{Quotient Set}
  \item \hyperref[def:canonical-surj]{Canonical Surjection}
  \item \hyperref[def:function]{Function}
  \item \hyperref[def:composition]{Composition}
\end{itemize}
\end{dependencies}

\begin{remark*}[Duality of perspectives]
This theorem establishes a bijection between equivalence relations on $A$
and partitions of $A$. The two perspectives---``same block'' (partition) and
``related'' (equivalence relation)---are interchangeable and each is more
natural in different contexts.
\end{remark*}

\begin{example*}[Extremal equivalence relations]
On any set $A$:
\begin{enumerate}
\item The \emph{equality relation} $(a,b) \in R \iff a=b$ gives singleton
classes $[a] = \{a\}$ and the finest partition of $A$.
\item The \emph{universal relation} $R = A \times A$ gives $[a] = A$ for all
$a$, and the coarsest partition (one block).
\end{enumerate}
All other equivalence relations on $A$ lie strictly between these extremes.
\end{example*}

% ---------------------------------------------------------
% Universal Property of Quotient Maps
% ---------------------------------------------------------

\begin{theorembox}{Theorem (Universal Property of the Quotient Map)}
\begin{theorem}[Universal Property of the Quotient Map]
\label{thm:quotient-universal}
Let $R$ be an equivalence relation on $A$, let $\pi : A \to A/R$ be the
canonical surjection, and let $f : A \to B$ be any function. Then $f$ is
\emph{constant on equivalence classes} --- meaning $(a,b) \in R$ implies
$f(a) = f(b)$ --- if and only if there exists a unique function
$\bar{f} : A/R \to B$ such that
\[
f = \bar{f} \circ \pi.
\]
When it exists, $\bar{f}$ is defined by $\bar{f}([a]) = f(a)$.
\hyperref[prf:quotient-universal]{Go to proof.}
\end{theorem}
\end{theorembox}

\begin{remark*}[Forward reference]
This universal property depends on functions, surjectivity, and composition.
It is stated here to complete the quotient-set picture, but its dependency
ceiling is intentionally forward-looking.
\end{remark*}

\begin{remark*}[Standard quantified statement]
$\bigl(\forall a,b \in A,\; (a,b)\in R \to f(a) = f(b)\bigr) \;\leftrightarrow\; \bigl(\exists!\, \bar{f} : A/R \to B,\; \forall a \in A,\; \bar{f}([a]) = f(a)\bigr)$.
\end{remark*}

\begin{dependencies}
\begin{itemize}
  \item \hyperref[def:equivalence-rel]{Equivalence Relation}
  \item \hyperref[def:canonical-surj]{Canonical Surjection}
  \item \hyperref[def:quotient-set]{Quotient Set}
  \item \hyperref[def:function]{Function}
  \item \hyperref[def:composition]{Composition}
\end{itemize}
\end{dependencies}

\begin{remark*}[Why this is called the universal property]
The condition $f = \bar{f} \circ \pi$ says the diagram commutes: going
from $A$ to $B$ directly via $f$ gives the same result as going first to
$A/R$ via $\pi$ and then to $B$ via $\bar{f}$. The word ``universal''
refers to the fact that $\pi : A \to A/R$ is the \emph{unique} quotient with
this property: any other map that ``respects the equivalence'' factors through
$\pi$.

The uniqueness of $\bar{f}$ is forced by $\pi$ being surjective: every
element of $A/R$ is $[a]$ for some $a$, so $\bar{f}([a]) = f(a)$ is the
only possible definition.
\end{remark*}

\begin{remark*}[Well-definedness is the key verification]
The heart of the theorem is well-definedness: the formula $\bar{f}([a]) =
f(a)$ could be ambiguous if $[a] = [b]$ but $f(a) \neq f(b)$. The hypothesis
that $f$ is constant on equivalence classes is precisely what prevents this:
$(a,b) \in R$ (i.e., $[a] = [b]$) implies $f(a) = f(b)$, so the output
does not depend on which representative of the class is chosen.

This ``well-definedness check'' recurs throughout mathematics. Whenever a
function is defined on equivalence classes by a formula involving
representatives, the first obligation is to verify that the formula gives
the same answer for all representatives. The universal property tells exactly
what condition on $f$ makes this work.
\end{remark*}

\begin{remark*}[Where this appears downstream]
This theorem is not an abstract curiosity. It is the engine behind:
\begin{itemize}
\item \emph{Integer construction}: the integers $\mathbb{Z}$ are defined as
equivalence classes of pairs of natural numbers, and arithmetic operations
are defined by formulas on representatives. The universal property verifies
the operations are well-defined.
\item \emph{Rational construction}: $\mathbb{Q}$ is defined as equivalence
classes of pairs $(p,q)$ with $q \neq 0$, and addition/multiplication on
representatives must be checked.
\item \emph{Quotient groups in algebra}: the quotient group $G/N$ is defined
by the same factorization, and the quotient homomorphism is $\pi$.
\item \emph{Quotient topology}: the quotient topological space is defined as
the finest topology on $A/R$ making $\pi$ continuous --- a condition stated
precisely via the universal property.
\end{itemize}
\end{remark*}

